\documentclass[10pt,a4paper]{article}

\usepackage[margin=1.8cm]{geometry}
\usepackage[T1]{fontenc}
\usepackage[utf8]{inputenc}
\usepackage{lmodern}
\usepackage{xcolor}
\usepackage{enumitem}
\usepackage{titlesec}
\usepackage[hidelinks]{hyperref}
\usepackage{microtype}

\definecolor{linkcolor}{HTML}{356F7C}
\hypersetup{colorlinks=true,urlcolor=linkcolor,linkcolor=linkcolor}
\setlength{\parindent}{0pt}
\setlength{\parskip}{4pt}
\setlist[itemize]{leftmargin=1.2em,itemsep=1pt,topsep=3pt}
\setlist[enumerate]{leftmargin=1.5em,itemsep=5pt,topsep=3pt}
\titleformat{\section}{\large\bfseries}{}{0pt}{}[\vspace{2pt}\hrule]
\titlespacing*{\section}{0pt}{13pt}{6pt}
\pagestyle{empty}

\begin{document}

{\LARGE\bfseries Eduardo Sampaio Pimenta}\par
\vspace{3pt}
Mathematician --- Probability Theory and Stochastic Processes\par
\href{mailto:sampaieduardo@gmail.com}{sampaieduardo@gmail.com}\quad
\href{https://lattes.cnpq.br/6512284414514021}{Lattes}\quad
\href{https://scholar.google.com/citations?hl=en\&user=SJJP-ngAAAAJ}{Google Scholar}\par

\section*{Research interests}
Probability theory, stochastic processes, scaling limits, invariance principles, boundary random walks, Feller processes, dynamic percolation and evolving random environments.

\section*{Research group}
\textbf{PROFE --- Grupo de Pesquisa em Probabilidade e F\'isica Estat\'istica}\par
Member of the research line
\href{http://dgp.cnpq.br/dgp/espelholinha/790496658105085770869}{\emph{Interacting Particle Systems}},
registered in the CNPq Directory of Research Groups.

\section*{Academic experience and education}
\begin{itemize}
  \item \textbf{Postdoctoral research}, Federal University of Minas Gerais (UFMG),Belo Horizonte, Brazil, 2026--present.
        Host/Supervisor: Remy de Paiva Sanchis.
  \item \textbf{Postdoctoral research}, University of Augsburg, Augsburg, Germany, 2025--2026.
        Host/Supervisor: Markus Heydenreich.
  \item \textbf{PhD in Mathematics}, Federal University of Bahia (UFBA), Salvador, Brazil, 2021--2025.
        Advisors: Tertuliano Franco and Dirk Erhard.
  \item \textbf{MSc in Mathematics}, Federal University of Bahia (UFBA), Salvador, Brazil, 2019--2021.
        Advisor: Tertuliano Franco.
  \item \textbf{BSc in Mathematics}, Federal University of Minas Gerais (UFMG), Belo Horizonte, Brazil, 2010--2017.
\end{itemize}

\section*{Publications and preprints}
\begin{enumerate}
  \item D. Erhard, T. Franco, M. Jara and E. Pimenta. ``A Functional Central Limit Theorem for the General Brownian Motion on the Half-Line.'' Accepted, forthcoming in \emph{Bernoulli}. \href{https://www.bernoullisociety.org/index.php/publications/bernoulli-journal/bernoulli-journal-papers}{Journal page}.
  \item J. C. Arroyave, E. Barros and E. Pimenta. ``Scaling Limit of Boundary Random Walks: A Martingale Problem Approach.'' Preprint, 2025. \href{https://arxiv.org/abs/2507.10528}{arXiv:2507.10528}.
\end{enumerate}

\section*{Theses}
\begin{itemize}
  \item \textbf{Functional Central Limit Theorem and Berry--Esseen Estimates for Non-Homogeneous Random Walks}. Doctoral thesis, UFBA, 2025. Advisor 1: Tertuliano Franco; Advisor 2: Dirk Erhard.
  \item \textbf{Construction of the Dawson--Watanabe Process}. Master's dissertation, UFBA, 2021. Advisor: Tertuliano Franco.
\end{itemize}

\section*{Selected talks and presentations}
\begin{itemize}
  \item ``Scaling Limit of Boundary Random Walks: A Martingale Problem Approach.'' XVII CLAPEM, Montevideo 2026.
  \item ``A Functional Central Limit Theorem and Weak Berry--Esseen Estimates for Non-Homogeneous Random Walks.'' IHES Summer School --- Statistical Aspects of Nonlinear Physics, Paris 2025.
  \item ``Scaling Limits of Boundary Random Walks to Feller-Type Processes on the Half-Line.'' Augsburg 2025.
  \item ``A Functional CLT for the General Brownian Motion on the Half-Line.'' IX Encontro da pos-graduacao em matematica da UFBA, Salvador 2024.
  \item ``Berry--Esseen Estimates Based on Generators with Applications.'' VIII Encontro da Pós-graduação da UFBA, Salvador 2023.
\end{itemize}

\vfill
{\small\color{gray}Last updated: \today. This is a preliminary version to be completed from the full academic CV.}

\end{document}
